\section{Introduction}\label{sec:intro}

A volatility forecast is rarely consumed as a point estimate. Position
limits, margin requirements, and Value-at-Risk (VaR) all depend on an
uncertainty statement attached to the forecast: an interval, a quantile, a
band the user is told to trust at a stated level. Online conformal methods
\citep{gibbs2021adaptive, zaffran2022adaptive, bhatnagar2023improved} have
made such statements available for any forecasting model, with
distribution-free guarantees that hold under distribution shift. Those
guarantees, however, are \emph{marginal}: correct on average over time.

This paper documents, at scale, that marginal validity conceals a
systematic conditional failure precisely where uncertainty statements
matter most. On a survivorship-bias-free panel of the 100 largest US
stocks over 2005--2025 (519{,}843 stock-days), adaptive conformal intervals
for next-day realized volatility achieve their promised 90\% coverage
almost exactly, yet cover only 82.9\% of stress-regime days, while
over-covering calm ones (Figure~\ref{fig:coverage}). The upper tail, the
direction that maps to risk-management losses, covers 90.1\% in stress when
95\% is promised. Validity on average is purchased by failing during
crises: the 2020 COVID crash (78.6\% coverage), the 2025 tariff shock
(76.5\%), and every other major volatility episode in two decades. The
same disease afflicts a zero-shot time-series foundation model: the native
prediction band of Chronos \citep{ansari2024chronos} is not merely too
narrow in stress but points the wrong way, missing on the upside 14.6\% of
the time against a 10\% budget while missing only 5.7\% on the downside.

We then show that most of this conditional deficit can be removed, and
we are precise about the part that cannot. Our calibration layer wraps any
point forecaster and combines three ingredients, each of which we isolate
in ablations. First, \emph{cross-sectional pooling}: per-regime conformal
thresholds are updated by every stock in the panel, so a rare regime
accrues calibration signal at the panel rate rather than the single-stock
rate ($\sim$100$\times$ more updates). Second, \emph{regime conditioning}:
thresholds are maintained per market regime, estimated online from public
state variables (transparent VIX-percentile bins in the headline
configuration; a soft-membership HMM variant is evaluated as robustness).
Third, \emph{self-tuning adaptation}: per-regime learning rates are
selected online by an expert-aggregation scheme with an issued-interval
corrector, removing the layer's hand-tuned learning rates. The result on the same
panel: stress coverage rises from 82.9\% to 88.3\% with other regimes flat
at nominal, average interval score statistically indistinguishable from
the best baselines (Model Confidence Set), and stress-day 99\% VaR breach
rates of 1.8\% against 3.1--6.7\% for GARCH-$t$, CAViaR, filtered
historical simulation, and adaptive conformal baselines. The identical
layer, applied around the Chronos median, produces bands within 0.2
percentage points of nominal in every regime.

The deficit that remains is as informative as the repair. Coverage on the
\emph{second day} after a spike enters the stress regime stays near
72--76\% under every conditioning scheme we tried: finer regime
partitions, freshness subgroups, forward-looking option-implied signals
(VIX9D), and even a deliberately leaky ``oracle'' calibrator given
full-sample smoothed regime memberships
(Figure~\ref{fig:dse}). Day-2 under-coverage is not a regime-detection
failure; it is a shift in the score distribution faster than any
bounded-step tracker can follow, and option prices do not anticipate it
either. We report this as an honest boundary of the approach, with a
finite-sample explanation.

Three framing points delimit the contribution. Point-forecast accuracy is
held fixed by design: our forecaster pool ties standard HAR and gradient
boosting benchmarks (Table~\ref{tab:e0}), and no accuracy claim is made;
the contribution is that the \emph{uncertainty statements} attached to
standard forecasts become conditionally trustworthy. The guarantees we
prove are deterministic, hold under arbitrary cross-sectional and serial
dependence, and condition on the calibrator's own filtered regime estimate
(the object a risk desk can actually compute), with the gap to latent
truth quantified empirically at 0.3 percentage points. And the evaluation
is designed against survivorship: the universe is rebuilt each January
from CRSP point-in-time market capitalizations, so Lehman Brothers,
Wachovia, and Washington Mutual are in the panel through their deaths.

\subsection*{Contributions}
\begin{enumerate}
\item \textbf{The phenomenon.} The first large-scale documentation that
  marginally valid online conformal intervals for volatility are
  systematically conditionally invalid, with the failure concentrated in
  stress regimes and in the direction risk management cares about
  (Section~\ref{sec:results}).
\item \textbf{The method.} A regime-conditional, panel-pooled, self-tuning
  online conformal layer with finite-sample per-regime coverage guarantees
  that hold for arbitrary dependence, plus a one-sided head that maps to
  VaR (Sections~\ref{sec:method}--\ref{sec:theory}).
\item \textbf{The decomposition.} Ablations attributing the repair to its
  parts: pooling contributes most of the average stress coverage; regime
  conditioning contributes transition coverage, conditional VaR balance,
  and the per-regime guarantee itself; adaptivity removes all tuning
  (Section~\ref{sec:ablations}).
\item \textbf{Foundation-model repair.} Evidence that a zero-shot TSFM's
  native uncertainty is directionally miscalibrated in crises, and that
  the same layer repairs it. This is timely, given that current TSFM
  evaluations in volatility forecasting consider point accuracy only
  \citep{brini2026forecasting}.
\item \textbf{An honest limit.} The day-2 onset window resists every
  conditioning scheme, including oracle regimes and option-implied
  signals; we characterize why (Section~\ref{sec:onset}).
\end{enumerate}
